% article_afghanistan_seismicity.tex
% Destructive crustal earthquakes beyond the Hindu Kush nest (Afghanistan).
% Submission version for Geofisica Internacional.
% Compile: pdflatex; bibtex; pdflatex; pdflatex.


\documentclass[11pt]{article}
\usepackage[utf8]{inputenc}
\usepackage[T1]{fontenc}
\usepackage{lmodern}
\usepackage{amsmath,amssymb}
\usepackage{graphicx}
\graphicspath{{figures/}{./}}   % MANDATORY; \includegraphics takes bare content names
\usepackage{booktabs}
\usepackage{array}
\usepackage{siunitx}
\usepackage{listings}
\usepackage{xcolor}
\usepackage{algorithm}
\usepackage{algpseudocode}
\usepackage[numbers,sort&compress]{natbib}
\usepackage{geometry}\geometry{margin=2.5cm}
\usepackage{authblk}
\usepackage[colorlinks=true,allcolors=blue]{hyperref}
\usepackage{cleveref}

\lstdefinestyle{py}{language=Python,basicstyle=\ttfamily\footnotesize,
  breaklines=true,frame=single,numbers=left,numberstyle=\tiny,
  keywordstyle=\color{blue},commentstyle=\color{gray},showstringspaces=false}
\lstdefinestyle{sh}{language=bash,basicstyle=\ttfamily\footnotesize,
  breaklines=true,frame=single,numbers=left,numberstyle=\tiny,
  keywordstyle=\color{blue},commentstyle=\color{gray},showstringspaces=false}
\lstset{style=py}

\title{Destructive Crustal Earthquakes Beyond the Hindu Kush Nest: Statistical
Seismicity Regimes, Strain and Fault Structure of Afghanistan}

\author[1]{Polina Lemenkova\thanks{Corresponding author:
\href{mailto:lemenkova@itu.edu.tr}{lemenkova@itu.edu.tr};
ORCID: \url{https://orcid.org/0000-0002-5759-1089}}}
\author[1]{Abdullah Can Z\"{u}lfikar\thanks{ORCID:
\url{https://orcid.org/0000-0001-6610-3334}}}
\affil[1]{Disaster and Emergency Management Department, Institute of Earthquake
Engineering and Disaster Management, Istanbul Technical University, Maslak,
34469 Istanbul, T\"{u}rkiye}
\date{}

\begin{document}
\maketitle

% ===========================================================================
\begin{abstract}
\noindent
Afghanistan lies at the junction of the India--Eurasia collision, the Chaman
transform boundary and the Makran subduction system, and it suffers recurrent
lethal earthquakes. Seismic hazard studies of the country have concentrated on the
intermediate-depth Hindu Kush nest. However, the destructive events of 2022, 2023
and 2025 were shallow crustal earthquakes in the plate-boundary zone and in the
slowly deforming interior. This study compares the statistical seismicity regimes of
the intermediate-depth, plate-boundary and interior crustal domains of Afghanistan.
The workflow applies declustering and unsupervised clustering to open earthquake
catalogues for 1900--2025. It then derives frequency--magnitude statistics,
focal-mechanism classes and Kostrov strain rates in Python, with maps produced in
the Generic Mapping Tools. The merged catalogue holds
13\,777 events of moment magnitude 4 and above, of which 45\,\% are deeper than
70\,km. Above $M_w$ 5.5, the b-value is 0.69 in the western interior, 0.97 in the
intermediate-depth zone and 1.33--1.38 in the north-eastern crust and the
plate-boundary zone. The crustal domains released 55\,\% of the seismic moment
since 1900. They also hosted every destructive earthquake of 2022--2025. The results provide the first reproducible, nationwide statistical
regionalisation of Afghan seismicity linked to the mapped fault structure.
They give the Afghanistan National Disaster Management Authority a basis for
hazard zoning that separates shallow crustal sources beneath populated valleys
from the deep Hindu Kush nest.
\end{abstract}

\noindent\textbf{Keywords:} Afghanistan; Hindu Kush; Chaman fault; seismicity;
b-value; Kostrov strain; focal mechanisms; GMT

% ===========================================================================
\section{Introduction}\label{sec:intro}

Earthquakes cause a disproportionate share of disaster fatalities in countries where
rapid population growth meets fragile building stock and limited preparedness
\cite{Bilham2009, Holzer2013, Ambraseys2011, Jaiswal2010}. Afghanistan is a severe
example of this combination. Its historical record documents repeated destruction of
towns and villages in the Hindu Kush, the Kabul basin and the Herat region, and
unreinforced adobe and stone dwellings remain the dominant housing type
\cite{Ambraseys2003, Akhundzadah2025}. Global risk models identify the continental interiors of the
Alpine--Himalayan belt, including Afghanistan, as regions of high earthquake risk
\cite{Silva2020, England2011}. The four destructive events of
2022--2025 confirmed this assessment. The Herat sequence of October 2023 alone caused
damage that required an international rapid damage estimation
\cite{WorldBank2024}.

The seismicity of Afghanistan is governed by the northward motion of the Indian and
Arabian plates relative to Eurasia. This convergence drives crustal shortening,
lateral extrusion and deep lithospheric subduction across south-central Asia
\cite{Molnar1975, Tapponnier1982, Jackson1984, Bird2003}. Plate models predict
India--Eurasia relative motion of about 40\,mm\,yr$^{-1}$ near Afghanistan, and
geodetic networks show how this motion is partitioned between strike-slip and
thrust structures \cite{DeMets1994, DeMets2010, Vernant2004, Mohadjer2010,
Szeliga2012}. The country itself is a mosaic of Gondwana-derived blocks accreted to
Eurasia along Mesozoic sutures that were later reactivated as the Herat and
Central Badakhshan fault systems \cite{Tapponnier1981, Doebrich2006}. To the south
and east, the Makran subduction zone, the Sulaiman fold belt and the Zagros
fold-and-thrust belt frame the Afghan blocks and record the same convergence in
contrasting structural styles \cite{Byrne1992, Walker2004, Lemenkova202015,
Talebian2004, Jackson1988, Lemenkova202135, Vernant2004, Wheeler2005,
Lemenkova202216, Ruleman2007, Szeliga2012}.

\begin{figure}[htbp]\centering
  \includegraphics[width=0.90\textwidth]{fig_tectonic_setting}
  \caption{Regional tectonic setting of Afghanistan between 55--80$^\circ$E and
  25--42$^\circ$N. Colour-shaded relief from the SRTM15+ grid is shown at full
  intensity inside Afghanistan and faded outside it. Purple lines are PB2002 plate
  boundaries and the dashed polygon is the PB2002 Persia--Tibet--Burma orogen.
  Quaternary faults are coloured by kinematics (red, thrust or reverse; blue,
  strike-slip; green, normal) and drawn with widths that increase with activity
  class. Magenta circles are earthquakes of magnitude 4 and above deeper than
  70\,km, 1900--2025 ($n = 6327$). Arrows give the velocities of the Indian and
  Arabian plates relative to Eurasia computed from the PB2002 Euler poles. The inset
  globe locates the map frame. Software used for plotting figure: GMT 6.6.0,
  GhostScript 10.06.0. Data: SRTM15+; PB2002; USGS OFR 2007-1103; USGS ComCat;
  GSHHG; DCW. Source: authors.}
  \label{fig:tectonic}
\end{figure}

The principal active structures of Afghanistan respond to this convergence in
distinct ways (\cref{fig:tectonic}). The left-lateral Chaman fault system
accommodates a large part of the relative motion between India and Eurasia, and
geodetic studies show that segments of it creep while others remain locked
\cite{UlHadi2013, Fattahi2016, Barnhart2017, Dalaison2021}. Beneath the Hindu Kush,
a steep slab of lithosphere produces one of the most active intermediate-depth
earthquake nests on Earth, with large events at 150--250\,km depth
\cite{Pegler1998, Pavlis2000, Sippl2013, Schneider2013}. The origin of this slab,
its possible detachment and the mechanics of its earthquakes have been investigated
through seismic tomography, moment-tensor analysis and thermal modelling
\cite{Lister2008, Kufner2016, Kufner2017, Zhan2016, Poli2016, Molnar2019, Lou2009}.
In contrast, the right-lateral Herat fault and the fold-and-thrust belts along the
northern margin of the Afghan blocks deform slowly and have produced few
instrumentally recorded earthquakes \cite{Wheeler2005, Ruleman2007}.

Earlier hazard assessments of Afghanistan relied on these catalogues and fault maps.
The national probabilistic hazard map combined a seismotectonic synthesis, a
regional earthquake catalogue and a database of Quaternary faults
\cite{Boyd2007, Dewey2006, Wheeler2005, Ruleman2007}. Global programmes have since
incorporated Afghanistan into harmonised hazard and active-fault models
\cite{Giardini1999, Styron2020, Silva2020}, and recent reviews have emphasised the
growing toll of Afghan earthquakes \cite{Akhundzadah2025}. Source studies of the
recent destructive events have shown that the 2023 Herat sequence ruptured
previously unmapped north-dipping blind thrusts at 8--14\,km depth in five
successive mainshocks \cite{Zhao2025, Fu2025, Xu2025, Zhong2026, Toker2026}. The
2022 Paktika earthquake produced strong ground motion and surface deformation along
the plate-boundary zone \cite{Jahed2024}, and the 2015 Hindu Kush event damaged
buildings across northern Pakistan despite its 231\,km focal depth
\cite{Ismail2016}.

These studies leave a quantitative gap. Seismicity in Afghanistan has been analysed
either event by event or within continental-scale compilations. No study has yet
compared the statistical behaviour of its distinct seismotectonic domains with a
single, reproducible workflow. Catalogue-wide statistics are dominated by
intermediate-depth events, mostly beneath the Hindu Kush, which supply nearly half
of all recorded earthquakes but rarely produce lethal shaking. The shallow crustal domains that generated every lethal
earthquake of 2022--2025 are, in contrast, sparsely sampled by the instrumental
record. The historical record of damaging earthquakes shows this contrast clearly
(\cref{fig:damaging}). The NOAA database lists 126 earthquakes with reported deaths
in the region since 1900, and 45 of them affected Afghanistan \cite{NOAA_SigEQ}.
Lethal shaking in Afghanistan clusters in the north-east and along the eastern and
western margins. Reported lethal events become markedly more frequent from the
1990s onwards, which partly reflects improved reporting. Experience in slowly deforming continental interiors elsewhere shows that
short catalogues can underrepresent the hazard of such regions, so that hazard maps
built on them fail where they matter most \cite{Stein2012, Liu2016, Calais2016,
England2011}.

Domain-resolved statistics are needed to direct risk reduction in Afghanistan,
where national disaster management depends on hazard zonation that reflects the
types of earthquake each region produces. Machine-learning tools meet this need.
They partition seismic catalogues objectively and have already improved the
characterisation of earthquake clusters and seismic hazard in other collision zones \cite{Bergen2019, Kong2018, Lemenkova202606, Mousavi2022,
Zaliapin2016}. Their combination with classical statistical seismology and
open-source cartography offers a transparent route to a nationwide comparison that
partner institutions in Afghanistan can reproduce and extend.

This study aims to quantify and compare the statistical seismicity regimes of the
intermediate-depth, plate-boundary and interior crustal domains of Afghanistan using
open data only. The contributions are:
\begin{itemize}
  \item a merged, magnitude-homogenised and declustered open catalogue of
        Afghanistan and its margins for 1900--2025;
  \item a data-driven delineation of seismotectonic domains by unsupervised
        density-based clustering of hypocentres;
  \item domain-resolved completeness magnitudes, b-values, depth distributions and
        recurrence estimates with bootstrap uncertainties;
  \item focal-mechanism classification and seismic strain-rate fields from Kostrov
        summation of moment tensors;
  \item a quantitative comparison of mapped fault structure with the spatial
        distribution of seismicity;
  \item a series of reproducible maps produced with the Generic Mapping Tools and
        openly released scripts.
\end{itemize}
\Cref{sec:studyarea} describes the tectonic setting, \cref{sec:data} the open
datasets and \cref{sec:methods} the methodology. \Cref{sec:results} reports the
results, which are interpreted in \cref{sec:discussion} and summarised in
\cref{sec:conclusions}.

\begin{figure}[htbp]\centering
  \includegraphics[width=0.72\textwidth]{fig_damaging_events}
  \caption{Damaging earthquakes in Afghanistan and adjacent regions, 1900--2025.
  (a)~Epicentres of 205 significant earthquakes in 58--76\,$^\circ$E,
  27--41\,$^\circ$N on a shaded-relief base, sized by magnitude and coloured by
  reported fatalities. Labels give the year and death toll of the most lethal events
  and of the recent Afghan sequence. (b)~Fatalities per event for the 126
  earthquakes with reported deaths on a logarithmic axis. The 45 events that affected
  Afghanistan are shown as opaque outlined bars and the events in neighbouring
  countries as translucent bars. Software used for plotting figure: GMT 6.6.0.
  Data: NOAA NCEI/WDS Global Significant Earthquake Database \cite{NOAA_SigEQ};
  SRTM15+ relief. Source: authors.}
  \label{fig:damaging}
\end{figure}

% ===========================================================================
\section{Study Area and Tectonic Setting}\label{sec:studyarea}

Afghanistan covers about 652\,000\,km$^2$. It extends from 29.4 to 38.5 degrees
north and from 60.5 to 74.9 degrees east. Its relief rises from the Helmand and
Sistan basins in the south-west, at about 500\,m, to peaks above 7000\,m in the
eastern Hindu Kush and the Wakhan corridor (\cref{fig:geology}). The northern
plains of the Amu Darya basin form the foreland of the Afghan--Tajik depression,
whereas the south-eastern ranges pass into the Sulaiman and Kirthar fold belts of
Pakistan \cite{Tapponnier1981, Doebrich2006}. This framework defines
three settings: the high mountains of the north-east, the transpressional
plate-boundary zone of the east and the plateau and basins of the centre and west.

Geologically, the country comprises the Eurasian margin north of the Herat fault,
the Farah, Helmand and Kabul blocks to the south, and the Katawaz flysch basin
adjacent to the Chaman fault \cite{Tapponnier1981, Doebrich2006, Wheeler2005}. These
units were assembled along Mesozoic sutures, which were reactivated during the
Cenozoic as major strike-slip and thrust faults. The Herat fault extends for about
700\,km from the Iranian border to its junction with the Chaman fault north of
Kabul \cite{Zhao2025, Ruleman2007}. Foreland basins on the margins of the
collision system record the same convergence far from the plate boundary. The
Mesopotamian foredeep of Iraq is one example, shaped by sedimentary loading and
flexure \cite{Jackson1984, Talebian2004, Lemenkova202132, Vernant2004,
Walker2004}.

Historical earthquakes reveal damaging activity in all three settings. Large
shallow events struck Herat in AD\,847, 1102 and 1364 and the Bamyan region in
1950, and earthquakes along the Chaman and Kabul fault systems repeatedly damaged
Kabul and Jalalabad \cite{Ambraseys2003, Zhao2025}. Intermediate-depth events of
magnitude 7 and above beneath the Hindu Kush recur at intervals of years to decades
and are felt over much of south-central Asia \cite{Dewey2006, Pegler1998,
Zhan2016}. Having outlined this framework, the next section describes the open
datasets used to characterise it quantitatively.

\begin{figure}[htbp]\centering
  \includegraphics[width=0.92\textwidth]{fig_geology_faults}
  \caption{Geology and faults of Afghanistan. The 11\,202 polygons of the USGS
  1:850\,000 geological map are grouped into ten chronostratigraphic classes in the
  colours of the International Chronostratigraphic Chart and drawn semi-transparently
  over SRTM15+ shaded relief. Thin dark lines are the 14\,885 mapped faults of the same
  map, and thicker grey lines its 83 thrusts. Quaternary faults are coloured by
  kinematics: reverse in crimson, strike-slip in black and normal in blue. Dashed
  lines are provincial boundaries. Labels mark the principal cities, tectonic blocks
  and the Herat and Chaman faults. Software used for plotting figure: GMT 6.6.0,
  GDAL 3.12. Data: USGS geology and faults of Afghanistan \cite{Doebrich2006}; USGS
  Quaternary faults \cite{Ruleman2007}; geoBoundaries \cite{Runfola2020}; SRTM15+.
  Source: authors.}
  \label{fig:geology}
\end{figure}

% ===========================================================================
\section{Data and Materials}\label{sec:data}

All data used in this study are open access and were retrieved from public
repositories; no field data or restricted datasets were used (\cref{tab:data}).

\subsection{Earthquake Catalogues}
The primary catalogue is the Comprehensive Catalog of the United States Geological
Survey (ComCat). It was queried for events of magnitude 4 and above between 1900
and 2025 in the window 58--76$^\circ$E and 27--41$^\circ$N. ComCat merges contributions from
global and regional agencies and underpins the catalogues used for national hazard
models \cite{Mueller2018}. For the period before 1964 and for large events, the
ISC-GEM catalogue provides relocated hypocentres and moment magnitudes derived with
uniform procedures \cite{Storchak2013, DiGiacomo2018}. Its solutions were obtained
through ComCat under the contributor codes \texttt{iscgem} and \texttt{iscgemsup},
which yielded 2078 events in the study window. The bulletin of the
International Seismological Centre and the relocated catalogue of Engdahl and
colleagues are used to check hypocentral depths, which are critical for separating
crustal from intermediate-depth events \cite{Bondar2011, Engdahl1998}.

\subsection{Focal Mechanisms}
Centroid moment tensors for 1976--2025 are taken from the Global Centroid Moment
Tensor (GCMT) project, which provides full moment tensors for earthquakes of
moment magnitude above about 5 worldwide \cite{Dziewonski1981, Ekstrom2012}. These
solutions supply the faulting style, principal-axis orientations and scalar moments
used in the mechanism and strain analyses.

\subsection{Geological, Fault and Topographic Data}
Geological units and mapped faults are taken from the digital geological and mineral
resource map of Afghanistan at 1:850\,000, which contains 11\,202 lithological
polygons and 14\,968 fault segments \cite{Doebrich2006}. Fault activity and
kinematics are taken from the database of probable and possible Quaternary faults,
which classifies 29\,619 segments into thrust, strike-slip and normal types and
three activity classes \cite{Ruleman2007}. The GEM Global Active Faults Database
provides a regional context beyond the national border \cite{Styron2020}. Plate
boundaries and orogen outlines follow the PB2002 model \cite{Bird2003}. Topography
is represented by the SRTM15+ grid at 15 and 30 arc-second resolution
\cite{Tozer2019}. Coastlines, rivers and political borders come from the
Global Self-consistent, Hierarchical, High-resolution Geography database (GSHHG)
\cite{Wessel1996}. Administrative boundaries are taken from geoBoundaries
\cite{Runfola2020}.

\subsection{Damage and Fatality Records}
Fatality and damage figures for historical and recent earthquakes are compiled from
the published Afghan earthquake record and from post-event reports
\cite{Ambraseys2003, WorldBank2024}. They serve only to identify destructive events
and are not used in the statistical analysis.

\begin{table}[htbp]\centering
  \caption{Open-access datasets used in the study, with provider, coverage and
  role in the analysis.}
  \label{tab:data}
  \small
  \begin{tabular}{p{3.4cm}p{3.6cm}p{3.0cm}p{4.2cm}}
    \toprule
    Dataset & Provider / version & Coverage & Role \\
    \midrule
    ComCat earthquake catalogue & USGS, accessed Nov 2026 & 1900--2025, M$\geq$4 & Primary catalogue \\
    ISC-GEM catalogue & ISC, via ComCat & 1905--2021, 2078 events & Early and large events, $M_w$ \\
    ISC Bulletin, EHB & ISC & 1964 onwards & Depth control \\
    Global CMT & GCMT project & 1976--2025 & Focal mechanisms, moments \\
    Geology and faults & USGS OFR 2006-1038 & 1:850\,000 & Lithology, fault density \\
    Quaternary faults & USGS OFR 2007-1103 & 29\,619 segments & Fault kinematics, activity \\
    PB2002 & Bird (2003) & Global & Plate boundaries \\
    SRTM15+ & Tozer et al. (2019) & 15 and 30 arc-sec & Relief basemaps \\
    GSHHG, geoBoundaries & Wessel and Smith; Runfola et al. & Global & Coasts, borders \\
    \bottomrule
  \end{tabular}
\end{table}

The assembled ComCat catalogue is summarised in \cref{fig:catalogue}, which shows
the spatial separation between the dense intermediate-depth cluster beneath the
Hindu Kush and the diffuse crustal seismicity elsewhere.

\begin{figure}[htbp]\centering
  \includegraphics[width=0.78\textwidth]{fig_catalogue_overview}
  \caption{Instrumental seismicity of Afghanistan and its surroundings,
  1900--2025. Circles mark 13\,779 earthquakes of magnitude 4 and above from the
  USGS ComCat catalogue. They are coloured by hypocentral depth from 0 to 300\,km
  and sized by magnitude, on a grey shaded-relief base from SRTM15+. Thin dark lines are
  Quaternary faults and the bold line is the national border. Stars mark the
  destructive earthquakes of Paktika (2022), Herat (2023), Kunar (2025) and Kholm
  (2025), labelled with catalogue magnitudes. Panel (a) shows the depth histogram
  in 10\,km bins and panel (b) the magnitude histogram in 0.1-unit bins, both on a
  logarithmic count axis. Software used for plotting figure: GMT 6.6.0,
  GhostScript 10.06.0. Data: USGS ComCat; USGS OFR 2007-1103; SRTM15+; GSHHG; DCW.
  Source: authors.}
  \label{fig:catalogue}
\end{figure}

These data are processed in five stages, described next.

% ===========================================================================
\section{Methodology}\label{sec:methods}

\subsection{Methodological Framework}
The workflow comprises five stages (\cref{fig:workflow}). The first stage merges the
catalogues and converts all magnitudes to moment magnitude. The second removes
dependent events, and the third delineates seismotectonic domains by unsupervised
clustering. The fourth stage computes completeness, frequency--magnitude, depth,
recurrence, mechanism and strain statistics for each domain. The fifth relates
seismicity to the mapped fault structure and produces the maps. All computations
use Python with NumPy, SciPy, pandas and scikit-learn
\cite{Harris2020, Virtanen2020, McKinney2010, Pedregosa2011}; graphs are drawn with
Matplotlib \cite{Hunter2007} and maps with the Generic Mapping Tools (GMT) version
6.6 \cite{Wessel2019}.

\begin{figure}[htbp]\centering
  \includegraphics[width=0.95\textwidth]{fig_workflow}
  \caption{Processing workflow. Open datasets (left) enter five stages: (1)
  catalogue merging and magnitude homogenisation; (2) declustering; (3) unsupervised
  delineation of seismotectonic domains; (4) statistical seismology, focal mechanisms
  and strain; and (5) fault-structure metrics and cartography. The main product of
  each stage is shown on the right, with the event counts of the final analysis.
  Software used for plotting figure: Matplotlib 3.10. Source: authors.}
  \label{fig:workflow}
\end{figure}

\subsection{Catalogue Homogenisation}
The ISC-GEM solutions distributed through ComCat under the contributor codes
\texttt{iscgem} and \texttt{iscgemsup} are merged with the ComCat catalogue. A
ComCat event is replaced by the ISC-GEM solution when the two origins differ by no
more than 16\,s and 50\,km. A second pass removes residual duplicates within the
same tolerances. It retains the ISC-GEM solution or, between two ComCat entries,
the larger event. Moment magnitudes of all types are retained. Body-wave ($m_b$) and
surface-wave ($M_S$) magnitudes are converted to moment magnitude with the global
relations of Scordilis \cite{Scordilis2006}. The conversion is checked against a
regional general orthogonal regression of matched $m_b$--$M_w$ pairs, with an
error-variance ratio of 0.2 \cite{Das2011}. Moment magnitude is linked to scalar
seismic moment through the relation of Hanks and Kanamori \cite{Hanks1979,
Kanamori1977}:
\begin{equation}
  M_0 = 10^{1.5 M_w + 9.1}, \label{eq:moment}
\end{equation}
where $M_0$ is the scalar seismic moment (N\,m) and $M_w$ is the moment magnitude.
Moments computed with \cref{eq:moment} are summed to obtain cumulative moment
release, which is dominated by the largest events in each domain \cite{Kagan2002}.
Events at the default depths of 10, 33 and 35\,km are flagged, because their depths
reflect location procedures rather than source depth.

Converted magnitudes are used for mapping and for the delineation of domains but not
for the frequency--magnitude analysis. A linear conversion with a slope below one
compresses the magnitude range and changes the apparent b-value by the inverse of
that slope, so that b-values from converted magnitudes depend on the conversion
chosen. The frequency--magnitude statistics therefore use the period 1964--2025 and
moment magnitudes of 5.5 and above, where nearly all events carry a directly
determined moment magnitude (\cref{sec:results}). The merging, duplicate removal and
magnitude conversion are implemented as shown in \cref{lst:catalogue}.

\begin{lstlisting}[caption={Merging of the ISC-GEM and ComCat records, removal of
duplicates within 16\,s and 50\,km and conversion of body-wave and surface-wave
magnitudes to $M_w$ (Python 3, pandas, NumPy; excerpt of analysis.py).},
label={lst:catalogue}]
tcc = cc["time"].values.astype("datetime64[ms]").astype(np.int64) / 1000.0
pairs, used = [], set()
for _, g in gem.iterrows():                      # ISC-GEM has priority
    j = np.where(np.abs(tcc - g["time"].value / 1e9) <= 16)[0]
    if len(j):
        dd = dist_km(g.latitude, g.longitude, cc.latitude.values[j],
                     cc.longitude.values[j])
        if dd.min() <= 50:
            k = j[np.argmin(dd)]
            used.add(k)
            pairs.append((cc.loc[k, "mag"], cc.loc[k, "magType"], g["mag"]))
gem["source"], rest = "ISC-GEM", cc.drop(index=list(used)).copy()
rest["source"] = "ComCat"
cat = pd.concat([gem, rest], ignore_index=True).sort_values("time")

def to_mw(m, t, rel="scordilis"):
    if t.startswith("mw"):
        return m
    if t in ("mb", "mblg", "mb_lg"):             # Scordilis (2006) or regional GOR
        return 0.85 * m + 1.03 if rel == "scordilis" else beta * m + alpha
    if t.startswith("ms"):
        return 0.67 * m + 2.07 if m <= 6.1 else 0.99 * m + 0.08
    return m

cat["mw"] = np.round([to_mw(m, t) for m, t in zip(cat.mag, cat.magType)], 2)
cat = cat[cat.mw >= 4.0].reset_index(drop=True)
cat["fixed_depth"] = cat.depth.round(1).isin([10.0, 33.0, 35.0])
\end{lstlisting}

\subsection{Declustering}
Aftershocks and foreshocks are removed because frequency--magnitude and recurrence
statistics assume independent events. Two methods are compared. The window method
of Gardner and Knopoff removes events within magnitude-dependent space--time windows
around each mainshock \cite{Gardner1974}. Here the spatial window is applied to
hypocentral rather than epicentral distance, so that crustal events above the Hindu
Kush slab are not attributed to intermediate-depth mainshocks. The cluster method of
Reasenberg \cite{Reasenberg1985} was not applied, because its interaction parameters
are calibrated for dense regional networks rather than teleseismic catalogues. The
nearest-neighbour method of Zaliapin and Ben-Zion uses the proximity of each event to
its nearest earlier neighbour \cite{Zaliapin2008, Zaliapin2013, Zaliapin2016}:
\begin{equation}
  \eta_{ij} = t_{ij}\, r_{ij}^{\,d_f}\, 10^{-b\, m_i}, \label{eq:nnd}
\end{equation}
where $\eta_{ij}$ is the proximity of event $j$ to an earlier event $i$, $t_{ij}$ is
the inter-event time (yr) and $r_{ij}$ is the hypocentral distance (km). The
exponent $d_f$ is the fractal dimension of hypocentres (1.6), $b$ is the
Gutenberg--Richter slope (1.0) and $m_i$ is the magnitude of the parent event. The
distribution of $\log_{10}\eta$ is fitted with a two-component Gaussian mixture, and
events below the point of equal membership probability are classified as clustered.
The nearest-neighbour background catalogue is used in all subsequent analyses, and
the Gardner--Knopoff catalogue serves as a sensitivity test.

\subsection{Data-Driven Delineation of Seismotectonic Domains}
Seismotectonic domains are delineated by density-based clustering rather than drawn
by hand. The algorithm HDBSCAN extends DBSCAN by building a hierarchy of clusters
across density levels and retaining the most persistent ones. It therefore requires
no fixed density threshold and labels sparse events as noise
\cite{Ester1996, Campello2013, McInnes2017}. Each background event is described by
its projected easting and northing (km) and its depth multiplied by three, a weight
that separates crustal from intermediate-depth seismicity in the same epicentral
area. The minimum cluster size is 80 events and the minimum number of samples is 20
(\cref{lst:clustering} in Appendix~\ref{secA1}).
Similar density-based approaches have been used to reconstruct fault networks from
seismicity \cite{Ouillon2008}.

The clusters are grouped into four domains by simple rules applied to their median
location. Clusters with a median depth of 70\,km or more form
the Hindu Kush--Pamir intermediate-depth zone (HKP). Crustal clusters with a median
longitude west of 64.5$^\circ$E form the western interior domain (WEST). Crustal
clusters south of 35.2$^\circ$N, or east of 72.5$^\circ$E and south of 36$^\circ$N,
form the plate-boundary domain (PBZ), and the remaining crustal clusters form the
north-eastern crustal domain (NEC). Every event of the catalogue, including those
labelled as noise or as clustered, is then assigned to a domain by a
nearest-neighbour classifier with 15 neighbours trained on the clustered events
\cite{Pedregosa2011}.


\subsection{Magnitude of Completeness and b-Value Estimation}
The magnitude of completeness $M_c$ is estimated by the maximum-curvature method with
a correction of $+0.2$ \cite{Wiemer2000, Woessner2005}. Its temporal evolution is
traced in moving windows of 400 events with bootstrap uncertainties. For the
frequency--magnitude analysis of moment magnitudes, $M_c$ is set to 5.5 for the
period 1964--2025 after checking the maximum-curvature estimate of the directly
determined moment magnitudes in each domain. The frequency--magnitude distribution
above $M_c$ follows the Gutenberg--Richter law \cite{Gutenberg1944}:
\begin{equation}
  \log_{10} N(\geq M) = a - b\,M, \label{eq:gr}
\end{equation}
where $N(\geq M)$ is the number of events with magnitude at least $M$, $a$ measures
seismic productivity and $b$ describes the relative frequency of small and large
events. The b-value is estimated by maximum likelihood \cite{Aki1965}:
\begin{equation}
  \hat{b} = \frac{\log_{10} e}{\bar{M} - \left(M_c - \Delta M/2\right)},
  \label{eq:aki}
\end{equation}
where $\bar{M}$ is the mean magnitude of events at or above $M_c$ and $\Delta M$ is
the magnitude bin width (0.1). Its standard error follows Shi and Bolt
\cite{Shi1982}:
\begin{equation}
  \sigma_b = 2.3\,\hat{b}^{2}\sqrt{\frac{\sum_{i=1}^{n}(M_i-\bar{M})^2}{n(n-1)}},
  \label{eq:shibolt}
\end{equation}
where $n$ is the number of events at or above $M_c$ and $M_i$ their magnitudes.
Confidence intervals are obtained from 1000 bootstrap resamples. Because the b-value
decreases with increasing differential stress, its variations between domains are
used as a proxy for their stress state \cite{Scholz1968, Schorlemmer2005,
Scholz2015}. Gridded b-values are computed from the nearest events to each grid node
(\cref{alg:bgrid}). Moment magnitudes are too sparse for spatial mapping, so the grid
uses the native body-wave magnitudes of the declustered catalogue for 1973--2025,
which make up 85\,\% of all events. Converting $m_b$ to $M_w$ with a slope of 0.85
would multiply every node value by the same factor and leave the spatial pattern
unchanged. The gridded values are therefore read as relative contrasts, and absolute
b-values are taken from the $M_w$ domain estimates. Crustal ($<70$\,km) and
intermediate-depth ($\geq 70$\,km) events are mapped separately at nodes spaced
0.25\textdegree\ apart, with $N = 150$, $R_{max} = 150$\,km and $n_{min} = 50$
(\cref{lst:bvalue}).

\begin{lstlisting}[caption={Completeness by maximum curvature, Aki--Utsu b-value and
Shi--Bolt standard error, applied per domain and per grid node (Python 3, NumPy).},
label={lst:bvalue}]
DM, NMIN = 0.1, 50                       # bin width, minimum events above Mc

def mc_maxc(m):
    bins = np.round(np.arange(m.min(), m.max() + DM, DM), 1)
    h = np.array([(m == b).sum() for b in bins])
    return float(bins[np.argmax(h)]) + 0.2  # maximum curvature + 0.2

def bfit(m, mc):
    x = m[m >= mc - 1e-9]
    n = len(x)
    if n < NMIN:
        return None
    mean = x.mean()
    b = np.log10(np.e) / (mean - (mc - DM / 2))       # Aki-Utsu
    sb = 2.3 * b**2 * np.sqrt(((x - mean)**2).sum() / (n * (n - 1)))  # Shi-Bolt
    return b, sb, n
\end{lstlisting}

\begin{algorithm}[htbp]
\caption{Gridded b-value mapping with nearest-$N$ sampling.}\label{alg:bgrid}
\begin{algorithmic}[1]
\State {select the nearest $N$ events within radius $R_{max}$ of each grid node}
\State {estimate $M_c$, $\hat{b}$ and $\sigma_b$ from the selected events}
\State {mask the node if fewer than $n_{min}$ events exceed $M_c$}
\State {write node values to a GMT grid for mapping}
\end{algorithmic}
\end{algorithm}

\subsection{Depth Distribution and Recurrence}
Hypocentral depths are summarised per domain by percentiles and kernel-density
estimates after excluding events at default depths. The bimodality of
intermediate-depth seismicity is quantified with a two-component Gaussian mixture
\cite{Pedregosa2011}. Recurrence is expressed as a Poisson probability. For each
domain, the probability of at least one event of magnitude $M$ or larger within
time $T$ is
\begin{equation}
  P(\geq M, T) = 1 - \exp\!\left(-10^{\,a - bM}\, T / T_{c}\right), \label{eq:poisson}
\end{equation}
where $T$ is the exposure time (yr) and $T_c = 62$\,yr is the duration of the
1964--2025 window used to estimate $a$ and $b$. The return period is the inverse of
the annual rate $10^{a-bM}/T_c$. Modelled return periods of $M_w \geq 7$ events are
compared with the observed numbers of such events in 1900--2025.

\subsection{Focal Mechanisms, Stress and Kostrov Strain}
Focal mechanisms are classified into thrust, strike-slip, normal and oblique types by
the plunges of their principal axes and displayed on ternary diagrams
\cite{Frohlich1992}. An event is a thrust when its T-axis plunge is at least
50\textdegree. It is a normal event when its P-axis plunge is at least 60\textdegree\
and a strike-slip event when its B-axis plunge is at least 60\textdegree. All other
events are oblique. Each Global CMT event is assigned to a domain by the same k-nearest-neighbour
classifier used for the catalogue (\cref{sec:methods}). The mean P axis of a domain is
the principal eigenvector of the summed orientation tensor of its P axes, which treats
the axes as undirected lines. The share of variance carried by this eigenvector
measures how tightly the axes cluster.
Seismic strain rates follow from the summation of moment tensors in each domain or
grid cell \cite{Kostrov1974, Jackson1988}:
\begin{equation}
  \dot{\varepsilon}_{ij} = \frac{1}{2\mu V T}\sum_{k=1}^{K} M_{ij}^{(k)},
  \label{eq:kostrov}
\end{equation}
where $\dot{\varepsilon}_{ij}$ is the strain-rate tensor (yr$^{-1}$) and $\mu$ is the
shear modulus ($3\times10^{10}$\,Pa). $V$ is the seismogenic volume (m$^3$), $T$ is
the catalogue duration (yr) and $M_{ij}^{(k)}$ is the moment tensor of event $k$
(N\,m). The summation uses the horizontal components of the Global CMT tensors for
1976--2025 ($T = 50$\,yr). Shallow events ($\le 70$\,km) are summed in
1\textdegree$\times$1\textdegree\ cells over a 30\,km seismogenic layer.
Intermediate-depth events (70--300\,km) of the Hindu Kush--Pamir zone are summed in
0.5\textdegree$\times$0.5\textdegree\ cells over a 230\,km layer. Only cells with
at least three events are retained (\cref{lst:kostrov}). Absolute rates scale inversely with the assumed
layer thickness, whereas the principal directions do not depend on it. The
eigenvectors of $\dot{\varepsilon}_{ij}$ give the principal shortening and extension
directions, which are compared with the geodetic velocity field
\cite{Mohadjer2010, Szeliga2012}. The horizontal scalar strain rate
$(\dot{\varepsilon}_{EE}^2 + \dot{\varepsilon}_{NN}^2 +
2\dot{\varepsilon}_{EN}^2)^{1/2}$ summarises the intensity of deformation in each
cell.

\begin{lstlisting}[caption={Kostrov summation of Global CMT moment tensors in one
grid cell, returning the principal horizontal strain rates, the azimuth of the
shortening axis and the scalar strain rate (Python 3, standard library).},label={lst:kostrov}]
def kostrov(events, lon0, lat0, d, H, mu=3.0e10, T=50.0, nmin=3):
    """Horizontal Kostrov strain rate (1/yr) in one d x d degree cell.
    events: (lon, lat, Mee, Mnn, Men) in N m; H: layer thickness in km."""
    s = [e for e in events if lon0 <= e[0] < lon0 + d and lat0 <= e[1] < lat0 + d]
    if len(s) < nmin:
        return None
    R = 6371.0e3
    area = R**2 * math.radians(d) * (math.sin(math.radians(lat0 + d))
                                     - math.sin(math.radians(lat0)))
    k = 1.0 / (2.0 * mu * area * H * 1e3 * T)
    ee, nn, en = (k * sum(e[q] for e in s) for q in (2, 3, 4))
    mean = 0.5 * (ee + nn)
    rad = math.hypot(0.5 * (ee - nn), en)
    e1, e2 = mean + rad, mean - rad                 # extension positive
    theta = 0.5 * math.degrees(math.atan2(2.0 * en, ee - nn))
    az_e2 = (180.0 - theta) % 180.0                 # azimuth of shortening axis
    scalar = math.sqrt(ee**2 + nn**2 + 2.0 * en**2)
    return e1, e2, az_e2, scalar
\end{lstlisting}

\subsection{Fault-Structure Metrics}
Fault density is gridded from the 14\,968 mapped fault segments:
\begin{equation}
  \rho_f(\mathbf{x}) = \frac{1}{A}\sum_{l \in A(\mathbf{x})} L_l, \label{eq:faultdens}
\end{equation}
Here $\rho_f$ is the fault density (km\,km$^{-2}$) and $A$ is the cell area (km$^2$).
$L_l$ is the length (km) of segment $l$ within the cell centred on $\mathbf{x}$; cells
are 0.25\textdegree\ wide. Fault strikes are summarised as length-weighted rose
diagrams in 10\textdegree\ sectors and compared with the nodal-plane strikes of
shallow Global CMT mechanisms. The distance from each shallow ($<40$\,km) epicentre
inside Afghanistan to the nearest Quaternary fault is compared with that of five
times as many uniformly random points inside the country. Faults are densified to
vertices at most 0.5\,km apart, and the two distance distributions are compared with
a two-sample Kolmogorov--Smirnov test to determine whether seismicity concentrates
near mapped active faults \cite{Ruleman2007, Styron2020}. The gridding and the
distance test are listed in \cref{lst:faults}.

\subsection{Cartographic Mapping}
All maps are produced with GMT scripts in Mercator projection, with shaded relief
from SRTM15+ and coastlines and borders from GSHHG. Depth and gridded statistics use
high-contrast continuous palettes (turbo and jet) over grey relief, each with a
labelled colour bar, and categorical classes combine distinct colours with distinct
symbols or outlines so that they remain separable in greyscale \cite{Crameri2020}. Hypocentres are coloured by depth and sized by magnitude, focal
mechanisms are drawn as lower-hemisphere projections, and gridded statistics are
rendered as continuous surfaces with masked low-data nodes. The scripts for the
seismicity map, the depth sections and the focal mechanisms are given in
\cref{lst:seismap,lst:sections,lst:meca}. The same scripted
approach has been applied by the authors to seismicity and geophysical mapping of
other active margins and rift systems, which ensures consistent cartographic
conventions across regions \cite{Wessel2019, Wessel1996, Lemenkova202218,
Tozer2019, Runfola2020, Lemenkova202205, Hunter2007, Harris2020, Lemenkova202215,
Virtanen2020, McKinney2010}. Subduction-related and rift-related settings mapped
with these tools provide structural analogues for the Makran and Hindu Kush
margins \cite{Talebian2004, Jackson1988, Lemenkova202127, Frohlich1992,
Ekstrom2012, Lemenkova202104, Dziewonski1981, Bird2003, Lemenkova202220,
Styron2020, Kagan2002}.

\subsection{Statistical Analysis and Uncertainty}
Differences in b-value between domains are expressed as the share of paired bootstrap
draws in which one domain exceeds the other \cite{Shi1982, Woessner2005}. All
b-values, return periods and recurrence probabilities are reported with bootstrap
95\,\% intervals. Sensitivity is assessed by repeating the b-value estimation for
the undeclustered and Gardner--Knopoff catalogues, for catalogues without
default-depth events, for completeness magnitudes of 5.3 and 5.8, and for the
converted catalogue with the Scordilis and regional conversions (\cref{tab:params}).
Depth distributions are compared with two-sample Kolmogorov--Smirnov tests
\cite{Virtanen2020}.

\begin{table}[htbp]\centering
  \caption{Parameters of catalogue merging, declustering, clustering, completeness
  and b-value estimation.}
  \label{tab:params}
  \small
  \begin{tabular}{p{3.6cm}p{5.6cm}p{4.2cm}}
    \toprule
    Step & Parameter & Value \\
    \midrule
    Merging & Time and distance tolerance & 16\,s, 50\,km \\
    Conversion & $m_b \rightarrow M_w$ (Scordilis) & $0.85\,m_b + 1.03$ \\
     & $M_S \rightarrow M_w$ (Scordilis) & $0.67\,M_S + 2.07$; $0.99\,M_S + 0.08$ \\
    Gardner--Knopoff & Distance metric & hypocentral \\
    Nearest neighbour & $d_f$, $b$, distance metric & 1.6, 1.0, hypocentral \\
     & Threshold & two-component Gaussian mixture \\
    HDBSCAN & Features & easting, northing, $3\times$depth (km) \\
     & Minimum cluster size, minimum samples & 80, 20 \\
    Domain assignment & Classifier & $k$-nearest neighbours, $k = 15$ \\
    Completeness & Method; moving window & maximum curvature $+0.2$; 400 events \\
    Frequency--magnitude & Window; $M_c$; bin width & 1964--2025; 5.5; 0.1 \\
    Uncertainty & Bootstrap resamples & 1000 \\
    Gridded b-value & Magnitude; window; node spacing & $m_b$; 1973--2025; 0.25\textdegree \\
     & $N$, $R_{max}$, $n_{min}$ & 150, 150\,km, 50 \\
    \bottomrule
  \end{tabular}
\end{table}

The results of these analyses are presented next, beginning with the properties of
the merged catalogue.

% ===========================================================================
\section{Results}\label{sec:results}

\subsection{Catalogue Composition and Completeness}
The merged catalogue contains 13\,777 earthquakes of moment magnitude 4 and above.
They span the years 1902 to 2025. Of these, 1845 solutions come from ISC-GEM and 11\,932 from
ComCat; 1834 ISC-GEM events matched a ComCat origin, and a further 13 residual
duplicates were removed. Converted magnitudes account for 11\,736 events (85.2\,\%),
nearly all of them body-wave magnitudes. Default depths occur for 4481 events
(32.5\,\%).

The 761 matched pairs of ComCat body-wave and ISC-GEM moment magnitudes span $m_b$
4.2--7.4. Their residuals with respect to the Scordilis conversion have a mean of
$-0.02$ and a standard deviation of 0.17 magnitude units. The regional orthogonal
regression gives $M_w = 0.92\,m_b + 0.66$, with a slope steeper than the global
value of 0.85.

The completeness magnitude of the converted catalogue decreases from 5.4 in the early
1970s to 5.0 after 1976 and falls intermittently to 4.7--4.8 after 2013
(\cref{fig:mc}a). Above $M_w$ 5.5, 96.7\,\% of the events since 1964 carry a directly
determined moment magnitude, and the maximum-curvature completeness of these events
is 5.3--5.4 in every domain. These values support $M_c = 5.5$ for the
frequency--magnitude analysis.

\begin{figure}[htbp]\centering
  	\includegraphics[width=0.98\textwidth]{fig_completeness}
  \caption{Completeness and declustering of the merged catalogue. (a) Magnitude of
  completeness of the converted catalogue from the maximum-curvature method with a
  $+0.2$ correction. It is computed in moving windows of 400 events advanced by 100
  events and plotted at the window centre. The shaded band spans the 16th--84th
  bootstrap percentiles, and the dashed line marks the completeness magnitude of
  5.5 used for the moment-magnitude analysis. (b) Cumulative number of events in the
  full catalogue, in the nearest-neighbour background catalogue and among the
  Gardner--Knopoff mainshocks, 1900--2025. Software used for plotting figure:
  Python 3.14, Matplotlib, pandas, NumPy. Data: USGS ComCat; ISC-GEM. Source:
  authors.}
  \label{fig:mc}
\end{figure}

The nearest-neighbour proximities form two modes at $\log_{10}\eta = -6.99$ and
$-4.15$, separated at $-5.77$. The background catalogue retains 11\,391 events
(82.7\,\%), whereas the Gardner--Knopoff windows retain 5064 mainshocks (36.8\,\%)
(\cref{fig:mc}b).

\subsection{Seismotectonic Domains}
HDBSCAN identified 14 clusters among the 11\,391 background events and labelled
5626 events as noise (\cref{fig:domains}). Six clusters lie at intermediate depth,
with median depths of 107--208\,km beneath the Hindu Kush and 129\,km beneath the
Pamir. Three crustal clusters lie in the north-east: the northern Pamir and Tajik
depression (902 events), north-eastern Afghanistan around 69.3$^\circ$E (231
events) and Badakhshan (392 events). In the last of these, 97\,\% of the events
have the default depth of 33\,km. Three clusters form the plate-boundary domain: the
Quetta--Sulaiman region (711 events), the Kabul--Kunar region (102 events) and the
2005 Kashmir rupture zone (158 events). The western interior comprises clusters in
eastern Iran (344 events) and at Gazli in Uzbekistan (122 events).

\begin{figure}[htbp]\centering
  \includegraphics[width=0.95\textwidth]{fig_domains}
  \caption{Seismotectonic domains of Afghanistan and its surroundings. Circles mark
  the 13\,777 earthquakes of the merged catalogue, coloured by domain (purple,
  Hindu Kush--Pamir intermediate-depth zone; green, north-eastern crust; orange,
  plate boundary; blue, western interior) and sized by moment magnitude. White
  numbered circles mark the medians of the 14 HDBSCAN clusters of the
  nearest-neighbour background catalogue. Thin dark lines are Quaternary faults,
  the bold line is the national border and stars mark the destructive earthquakes
  of 2022--2025. Software used for plotting figure: GMT 6.6.0, GhostScript 10.06.0,
  scikit-learn. Data: USGS ComCat; ISC-GEM; USGS OFR 2007-1103; SRTM15+; GSHHG;
  DCW. Source: authors.}
  \label{fig:domains}
\end{figure}

After assignment of all events, the four domains contain 6104 (HKP), 4294 (NEC),
2477 (PBZ) and 902 (WEST) events (\cref{tab:domains}). In HKP, 99.3\,\% of the events are deeper
than 70\,km, and the domain contains 97.8\,\% of all such events; each crustal
domain has fewer than 2.5\,\% of its events at these depths. The domains extend beyond the
national border: WEST includes eastern Iran, NEC the Tajik depression and the
Pamir, and PBZ the Sulaiman ranges and the Himalayan front of Pakistan. The 2022
Paktika and 2025 Kunar earthquakes fall in PBZ, the 2023 Herat sequence in WEST and
the 2025 Kholm earthquake in NEC.

\begin{table}[htbp]\centering
  \caption{Seismotectonic domains: number of events, events shallower than 70\,km,
  depth percentiles (10th, 50th and 90th, default depths excluded), median
  epicentre and share of the seismic moment released in 1900--2025.}
  \label{tab:domains}
  \small
  \begin{tabular}{lrrccc}
    \toprule
    Domain & Events & $<$70\,km & Depth (km) & Median epicentre & Moment share \\
    \midrule
    HKP & 6104 & 43 & 96 / 152 / 222 & 36.53$^\circ$N, 71.13$^\circ$E & 45.3\,\% \\
    NEC & 4294 & 4226 & 14 / 30 / 61 & 38.56$^\circ$N, 71.59$^\circ$E & 20.9\,\% \\
    PBZ & 2477 & 2423 & 14 / 25 / 59 & 32.86$^\circ$N, 70.25$^\circ$E & 18.9\,\% \\
    WEST & 902 & 888 & 10 / 15 / 59 & 34.37$^\circ$N, 60.40$^\circ$E & 14.9\,\% \\
    \bottomrule
  \end{tabular}
\end{table}

\subsection{Frequency--Magnitude Statistics}
Above $M_c = 5.5$ for 1964--2025, the declustered catalogue gives $b = 1.07 \pm 0.06$
for the whole region (95\,\% interval 0.96--1.21, $n = 364$). The domains differ
markedly (\cref{fig:fmd,tab:gr}). The western interior has the lowest value,
$b = 0.69 \pm 0.11$ (0.53--0.96), followed by the intermediate-depth zone with
$b = 0.97 \pm 0.07$ (0.84--1.12). The north-eastern crust and the plate boundary
have the highest values, $1.33 \pm 0.16$ and $1.38 \pm 0.19$. In paired bootstrap
draws, the b-value of WEST is lower than that of NEC and PBZ in more than 99.8\,\%
of cases and lower than that of HKP in 95.9\,\%. The b-value of HKP is lower than
those of NEC and PBZ in 99.4\,\% and 98.9\,\% of cases, whereas NEC and PBZ do not
differ significantly.

\begin{table}[htbp]\centering
  \caption{Gutenberg--Richter parameters of the declustered catalogue for
  $M_w \geq 5.5$, 1964--2025: number of events, share with directly determined
  moment magnitude, $a$-value for the 62-year window, b-value with the Shi--Bolt
  standard error, and bootstrap 95\,\% interval.}
  \label{tab:gr}
  \small
  \begin{tabular}{lccccc}
    \toprule
    Domain & $n$ & Native $M_w$ & $a$ & $b \pm \sigma_b$ & 95\,\% interval \\
    \midrule
    HKP & 158 & 96.8\,\% & 7.54 & $0.97 \pm 0.07$ & 0.84--1.12 \\
    NEC & 88 & 96.6\,\% & 9.27 & $1.33 \pm 0.16$ & 1.09--1.72 \\
    PBZ & 82 & 95.1\,\% & 9.48 & $1.38 \pm 0.19$ & 1.08--1.78 \\
    WEST & 36 & 100\,\% & 5.36 & $0.69 \pm 0.11$ & 0.53--0.96 \\
    All & 364 & 96.7\,\% & 8.44 & $1.07 \pm 0.06$ & 0.96--1.21 \\
    \bottomrule
  \end{tabular}
\end{table}

\begin{figure}[htbp]\centering
  \includegraphics[width=0.75\textwidth]{fig_fmd_domains}
  \caption{Frequency--magnitude distributions of the four seismotectonic domains
  for the declustered catalogue, 1964--2025. Markers give the annual rate of events
  at or above each moment magnitude in 0.1-unit bins. Circles denote the Hindu
  Kush--Pamir intermediate-depth zone, squares the north-eastern crust, triangles
  the plate boundary and diamonds the western interior. Lines are maximum-likelihood Gutenberg--Richter fits
  above the completeness magnitude of 5.5 (dotted line), with b-values and standard
  errors in the legend. Rates below 5.5 include converted magnitudes. Software used
  for plotting figure: Python 3.14, Matplotlib, NumPy. Data: USGS ComCat; ISC-GEM.
  Source: authors.}
  \label{fig:fmd}
\end{figure}

The gridded b-values resolve 756 crustal and 292 intermediate-depth nodes
(\cref{fig:bmap}). Their standard errors are small, with a median $\sigma_b$ of 0.15
for crustal and 0.14 for intermediate-depth nodes. The crustal map divides into three
broad zones. Low values, with medians of 1.62 and 1.54, cover the Hindu Kush and
Kabul block between 34.5\textdegree N and 38\textdegree N and the Himalayan front
of Kashmir. Intermediate values, with a median of 1.87, characterise the Chaman and
Sulaiman fold belts south of 32\textdegree N. High values, with a median of 2.29,
cover the Afghan--Tajik depression and the Tien Shan foreland north of
38\textdegree N, and a smaller high-b patch (median 2.20) lies between Kabul and
Jalalabad. The intermediate-depth map shows the strongest contrast within the Hindu
Kush slab. Nodes beneath its southern tip, between 34.5\textdegree N and
35.6\textdegree N, have the lowest values of the whole grid, with a median of 1.15.
The core of the nest has a median of 1.75 and its western flank a median of 2.04. The
Pamir slab farther north-east has a median of 1.71. Because absolute values depend on
the body-wave magnitude scale, these contrasts are interpreted relative to the median
of each layer, which is 1.77 for the crust and 1.71 for the intermediate depths
(\cref{fig:bmap}c).

\begin{figure}[htbp]\centering
  \includegraphics[width=0.62\textwidth]{fig_bvalue_map}
  \caption{Gridded b-values of Afghanistan and its margins from body-wave magnitudes of
  the declustered catalogue, 1973--2025. (a)~Crustal events ($<70$\,km);
  (b)~intermediate-depth events ($\geq 70$\,km) of the Hindu Kush--Pamir zone, whose
  extent is outlined by the dashed frame in (a). Each 0.25\textdegree\ node uses the
  nearest 150 events within 150\,km, and nodes with fewer than 50 events above $M_c$
  are left blank. Dashed contours show the Shi--Bolt standard error at intervals of
  0.1, and black lines in (a) are USGS Quaternary faults. (c)~Distribution of node
  b-values, with medians as dashed lines. Software used for plotting figure: GMT
  6.6.0. Data: merged ComCat and ISC-GEM catalogue; USGS Quaternary faults
  \cite{Ruleman2007}; SRTM15+. Source: authors.}
  \label{fig:bmap}
\end{figure}

\subsection{Depth Structure}
Of the 13\,777 events, 7580 (55.0\,\%) are shallower than 70\,km and 6197
(45.0\,\%) are at 70\,km or deeper; only eight events exceed 300\,km
(\cref{fig:catalogue}). The intermediate-depth population divides almost equally
between 70--150\,km (3068 events) and 150--300\,km (3121 events). In the HKP domain,
with default depths excluded, the 10th, 50th and 90th depth percentiles are 96, 152
and 222\,km (\cref{fig:depth}). A two-component Gaussian mixture places 53\,\% of
these hypocentres in a mode at $115 \pm 22$\,km and 47\,\% in a mode at
$204 \pm 26$\,km. The crustal domains have median depths of 30\,km (NEC), 25\,km
(PBZ) and 15\,km (WEST).

Four depth sections show how this structure varies across the region
(\cref{fig:sections}). They use the 9486 ComCat hypocentres with depths other than the
default values, within 50\,km of each line. Section A--A$'$ crosses the Hindu Kush
from north to south and contains 3699 hypocentres. Crustal events above 60\,km occur
along its whole length, whereas intermediate-depth events are confined to a zone about
150\,km wide. Within this zone, a sub-horizontal band at 80--130\,km extends southward
over about 150\,km, and a near-vertical column reaches 280\,km beneath its northern
part. The two parts correspond to the two depth modes of the HKP domain. Section
B--B$'$ crosses the Pamir from north-east to south-west. Its 411 hypocentres outline a
narrower zone that deepens south-westward from about 100\,km to 200\,km. The crust
along both sections hosts a separate shallow population, with most events above
50\,km. Sections C--C$'$ across the Chaman fault and D--D$'$ along the Herat fault
contain only 31 and 61 located hypocentres, all shallower than 60\,km, because most
events in these regions carry default depths.

\begin{figure}[htbp]\centering
 	 \includegraphics[width=0.75\textwidth]{fig_depth_histograms}
  \caption{Hypocentral depth distributions of the four seismotectonic domains for
  the nearest-neighbour background catalogue with default depths excluded. Solid
  lines are Gaussian kernel-density estimates normalised to their maxima, markers on
  the right show the median depth of each domain, and dashed lines are the two
  components of the Gaussian mixture fitted to the Hindu Kush--Pamir
  intermediate-depth zone. The number of events per domain is given in the legend.
  Software used for plotting figure: Python 3.14, Matplotlib, SciPy. Data: USGS
  ComCat; ISC-GEM. Source: authors.}
  \label{fig:depth}
\end{figure}

\begin{figure}[htbp]\centering
  \includegraphics[width=0.53\textwidth]{fig_depth_sections}
  \caption{Hypocentral depth sections. (a)~Index map of 9486 USGS ComCat hypocentres
  of magnitude 4 and above, 1900--2025, coloured by depth, with the four section
  lines; events at the default depths of 10, 33 and 35\,km are excluded. (b)~A--A$'$
  across the Hindu Kush; (c)~B--B$'$ across the Pamir; (d)~C--C$'$ across the Chaman
  fault; (e)~D--D$'$ along the Herat fault. Hypocentres within 50\,km of each line are
  coloured by depth and sized by magnitude, and SRTM15+ topography is shown above each
  section. Software used for plotting figure: GMT 6.6.0. Data: USGS ComCat; SRTM15+.
  Source: authors.}\label{fig:sections}
\end{figure}

\subsection{Large and Destructive Earthquakes}
The merged catalogue contains 29 earthquakes of $M_w$ 7 and above: 12 in HKP, 6 in
NEC, 6 in PBZ and 5 in WEST. The intermediate-depth events include those of 1909
($M_w$ 7.7), 1921 ($M_w$ 7.8), 1949, 1965, 1983, 2002 and 2015 ($M_w$ 7.5 at
220\,km). The crustal events include the 1935 Quetta ($M_w$ 7.6), 1956 central
Afghanistan ($M_w$ 7.3), 1997 Zirkuh ($M_w$ 7.2), 2005 Kashmir ($M_w$ 7.6) and 2013
Saravan ($M_w$ 7.7) earthquakes. The crustal domains together released 54.7\,\% of
the seismic moment of 1900--2025 (\cref{tab:domains}). The destructive events of
2022--2025 are all shallow (\cref{tab:events}). The Herat sequence comprised five
events of $M_w$ 5.9--6.3 at 8--14\,km depth between 7 and 15 October 2023, four of
them of $M_w$ 6.3.

\begin{table}[htbp]\centering
  \caption{Destructive earthquakes of 2022--2025 in the ComCat catalogue, with
  origin time (UTC), epicentre, depth and moment magnitude.}
  \label{tab:events}
  \small
  \begin{tabular}{lccccc}
    \toprule
    Event & Origin time (UTC) & Lat. ($^\circ$N) & Lon. ($^\circ$E) & Depth (km) & $M_w$ \\
    \midrule
    Paktika & 2022-06-21 20:54 & 33.02 & 69.46 & 4 & 6.0 \\
    Herat (1) & 2023-10-07 06:41 & 34.60 & 61.93 & 14 & 6.3 \\
    Herat (2) & 2023-10-07 07:12 & 34.54 & 61.89 & 8 & 6.3 \\
    Herat (3) & 2023-10-07 07:40 & 34.63 & 62.00 & 12 & 5.9 \\
    Herat (4) & 2023-10-11 00:41 & 34.54 & 62.05 & 8 & 6.3 \\
    Herat (5) & 2023-10-15 03:36 & 34.65 & 62.12 & 9 & 6.3 \\
    Kunar & 2025-08-31 19:17 & 34.71 & 70.81 & 8 & 6.0 \\
    Kholm & 2025-11-02 20:29 & 36.55 & 67.39 & 28 & 6.2 \\
    \bottomrule
  \end{tabular}
\end{table}

\subsection{Focal Mechanisms and Strain}
The Global CMT catalogue contains 886 moment tensors for 1976--2025 in the study
region, with $M_w$ 4.6--7.6 (\cref{fig:meca}). Of these, 505 have centroids at depths of
70\,km or less and 381 lie deeper. Thrust faulting dominates both groups. It accounts
for 48\% of the shallow mechanisms and 71\% of the intermediate-depth ones. Strike-slip
faulting is the second shallow style at 28\%, followed by oblique (18\%) and normal
faulting (6\%). Below 70\,km strike-slip and normal events almost vanish (5\% and 2\%),
and the remaining events are oblique reverse mechanisms (\cref{fig:meca}b, c).

The domains differ in their mixture of styles (\cref{tab:meca}). The intermediate-depth
zone is the most uniform: 72\% of its events are thrusts with near-horizontal P axes
trending N165\textdegree E. The plate-boundary domain combines thrusting on the
Sulaiman and Himalayan fronts (54\%) with strike-slip faulting along the Chaman system
(27\%) under a N--S P axis (N173\textdegree E). The north-eastern crust has the lowest share of
thrusts (43\%) and the highest share of normal events (8\%), with a NNW--SSE P axis
(N151\textdegree E). The western interior has the largest strike-slip share (32\%),
contains no normal mechanisms and is the only domain with a NE--SW P axis
(N041\textdegree E). The P axes of shallow events
cluster between N140\textdegree E and N010\textdegree E, and the fewest trend
WNW--ESE, between N060\textdegree E and N100\textdegree E (\cref{fig:meca}d). Ten intermediate-depth events fall in
crustal domains by the classifier and are excluded from their shallow statistics.

\begin{figure}[htbp]\centering
  \includegraphics[width=0.78\textwidth]{fig_focal_mechanisms}
  \caption{Focal mechanisms of Afghanistan and its margins, 1976--2025. (a)~Global
  CMT moment tensors as lower-hemisphere projections sized by magnitude on shaded
  relief with USGS Quaternary faults. The 505 mechanisms at depths of 70\,km or less
  are coloured by faulting class after \citet{Frohlich1992}; the 381
  intermediate-depth mechanisms are shown in grey beneath them. (b, c)~Ternary
  diagrams of the squared sines of the T-, B- and P-axis plunges for shallow and
  intermediate-depth events. (d)~Rose diagram of P-axis azimuths of shallow events in
  10\textdegree\ sectors, treated as axial data. Software used for plotting figure: GMT
  6.6.0. Data: Global CMT \cite{Dziewonski1981, Ekstrom2012}; USGS OFR 2007-1103;
  SRTM15+. Source: authors.}\label{fig:meca}
\end{figure}

\begin{table}[htbp]\centering
  \caption{Focal-mechanism styles by seismotectonic domain (Global CMT, 1976--2025).
  Counts are followed by percentages. Crustal domains include only events at depths of
  70\,km or less. The mean P axis is the axial mean of the P-axis orientations, given
  as azimuth/plunge; $C$ is the share of variance carried by the mean axis.}
  \label{tab:meca}
  \begin{tabular}{lrllllll}
    \toprule
    Domain & $N$ & Thrust & Strike-slip & Normal & Oblique & P axis & $C$ \\
    \midrule
    HKP & 373 & 268 (72) & 14 (4) & 5 (1) & 86 (23) & 165/11 & 0.63 \\
    NEC & 235 & 100 (43) & 65 (28) & 18 (8) & 52 (22) & 151/00 & 0.65 \\
    PBZ & 203 & 109 (54) & 55 (27) & 10 (5) & 29 (14) & 173/11 & 0.59 \\
    WEST & 65 & 33 (51) & 21 (32) & 0 (0) & 11 (17) & 041/06 & 0.72 \\
    All & 886 & 512 (58) & 159 (18) & 35 (4) & 180 (20) & 165/07 & 0.59 \\
    \bottomrule
  \end{tabular}
\end{table}

Kostrov summation resolves 57 shallow cells, which hold 440 of the 505 shallow
mechanisms, and 21 intermediate-depth cells (\cref{fig:strain}). Shallow scalar strain
rates span four orders of magnitude, from $7.6\times10^{-11}$ to
$2.9\times10^{-7}$\,yr$^{-1}$, with a median of $2.8\times10^{-9}$\,yr$^{-1}$
(\cref{fig:strain}c). The fastest cell contains the 2005 Kashmir rupture
($2.9\times10^{-7}$\,yr$^{-1}$). It is followed by one cell in the northern Pamir and
two cells along the Sistan suture zone of eastern Iran, all of $8\times10^{-8}$ to
$1.2\times10^{-7}$\,yr$^{-1}$. The five fastest cells carry 72\% of the summed shallow
rate, so that crustal deformation is concentrated along the margins of Afghanistan.
Shortening dominates 47 of the 57 shallow cells. The rate-weighted mean shortening axis
trends N021\textdegree E, while individual cells rotate from NE--SW in the west to
NNW--SSE in the Pamir (\cref{fig:strain}a). The interior of Afghanistan between
62\textdegree E and 66\textdegree E has no cell with three or more mechanisms.

The intermediate-depth field is more uniform. Its median rate is
$7.4\times10^{-10}$\,yr$^{-1}$, and the five cells of the Hindu Kush nest between
70\textdegree E and 71.5\textdegree E carry 96\% of the summed rate, with a maximum of
$1.4\times10^{-7}$\,yr$^{-1}$ (\cref{fig:strain}b). Shortening dominates 18 of the
21 cells, and the rate-weighted shortening axis trends N176\textdegree E with a high
resultant length of 0.83. The Hindu Kush slab therefore shortens almost exactly N--S,
in line with the convergence of India and Eurasia, whereas the Pamir cells farther
north-east deform at rates lower by one to three orders of magnitude.

The merged catalogue released a total seismic moment of
$6.7\times10^{21}$\,N\,m in 1900--2025, equivalent to a single event of $M_w$ 8.5
(\cref{fig:moment}). Release is strongly concentrated in space: the 10\,\% of the 624
cells with the highest moment account for 93\,\% of the total, and the two cells with
the highest values lie in the Hindu Kush nest (\cref{fig:moment}a). The HKP domain
released 45\,\% of the total, followed by NEC (21\,\%), PBZ (19\,\%) and WEST (15\,\%)
(\cref{tab:domains}). The timing of release differs between domains
(\cref{fig:moment}b). The HKP and NEC domains released most of their moment before
1960, and their curves have risen only gradually since then. In contrast, 40\,\% of
the PBZ moment and 56\,\% of the WEST moment were released after 2000, mainly in the
2005 Kashmir and 2013 Saravan earthquakes. The three largest events supply 43\,\% of
the HKP moment, 57\,\% of the PBZ moment and 65\,\% of the WEST moment, so the
long-term budget of the crustal domains rests on a few large earthquakes. The four
destructive earthquakes of 2022--2025, of $M_w$ 6.0--6.3, together released only
$8.6\times10^{18}$\,N\,m, or 0.13\,\% of the total. Their impact therefore reflects
shallow depth, proximity to settlements and building vulnerability rather than the
size of the moment release.

\begin{figure}[htbp]\centering
  \includegraphics[width=0.72\textwidth]{fig_strain_rate}
  \caption{Seismic strain rates from Kostrov summation of Global CMT moment tensors,
  1976--2025. (a)~Shallow events ($\le 70$\,km) in 1\textdegree\ cells over a
  30\,km layer; the dashed frame outlines panel (b). (b)~Intermediate-depth events
  (70--300\,km) of the Hindu Kush--Pamir zone in 0.5\textdegree\ cells over a
  230\,km layer. Colours show the logarithm of the horizontal scalar strain rate on
  the jet scale. Crosses show the principal horizontal strain rates normalised in each
  cell, with inward arrows for shortening and outward arrows for extension. Only cells
  with at least three events are shown. (c)~Distribution of cell strain rates.
  Software used for plotting figure: GMT 6.6.0. Data: Global CMT
  \cite{Dziewonski1981, Ekstrom2012}; SRTM15+. Source: authors.}\label{fig:strain}
\end{figure}

\begin{figure}[htbp]\centering
  \includegraphics[width=0.64\textwidth]{fig_moment_release}
  \caption{Seismic moment release in Afghanistan and its margins, 1900--2025.
  (a)~Cumulative seismic moment of all events of the merged catalogue in
  0.5\textdegree\ cells, computed as $M_0 = 10^{1.5 M_w + 9.1}$\,N\,m, on a logarithmic
  jet scale. Stars mark the 2022 Paktika, 2023 Herat, 2025 Kunar and 2025 Kholm
  earthquakes. (b)~Cumulative moment of each seismotectonic domain through time on a
  logarithmic axis. Numbered circles mark the largest events, and stars mark the
  2022--2025 destructive earthquakes on the curves of their domains. Software used for
  plotting figure: GMT 6.6.0. Data: merged USGS ComCat and ISC-GEM catalogue
  \cite{Storchak2013}; SRTM15+. Source: authors.}
  \label{fig:moment}
\end{figure}

\subsection{Fault Structure and Seismicity}
The mapped faults of Afghanistan have a total length of about 65\,700\,km
(\cref{fig:faultdens}a). Fault density reaches 0.56\,km\,km$^{-2}$ and exceeds
0.28\,km\,km$^{-2}$ in one-tenth of the 767 occupied cells. The densest belts follow
the Herat fault zone, the central Afghan ranges between 64\textdegree E and
68\textdegree E, and the Badakhshan region north of the Hindu Kush. Their strikes are
dominantly NE--SW, with a maximum at N050--060\textdegree E (\cref{fig:faultdens}b). The
Quaternary faults have a broader distribution with a maximum at
N090--100\textdegree E, which reflects the E--W Herat fault and the faults of the
north-eastern ranges (\cref{fig:faultdens}c). The nodal planes of the 103 shallow
Global CMT mechanisms inside Afghanistan strike mainly N--S to NNE--SSW, close to the
trend of the Chaman system and the eastern ranges (\cref{fig:faultdens}d).

Shallow seismicity clusters strongly around the Quaternary faults
(\cref{fig:faultdens}e). The 1351 epicentres shallower than 40\,km have a median
distance of 3.5\,km to the nearest Quaternary fault, compared with 11.3\,km for 6755
random points. Overall, 87\,\% of the epicentres lie within 10\,km of a Quaternary
fault and 99\,\% within 25\,km, whereas the corresponding shares for random points
are 46\,\% and 72\,\%. The two-sample Kolmogorov--Smirnov statistic is $D = 0.41$
($p < 10^{-100}$). Epicentres also fall in cells of higher fault density, with a
median of 0.099 against 0.064\,km\,km$^{-2}$ for random points. The interior of the
Helmand and Farah blocks, which has few mapped Quaternary faults, records few
instrumental earthquakes (\cref{fig:faultdens}a).

\begin{figure}[htbp]\centering
  \includegraphics[width=0.86\textwidth]{fig_fault_density}
  \caption{Fault structure and seismicity of Afghanistan. (a)~Density of the
  mapped faults of the USGS 1:850\,000 map (km\,km$^{-2}$) in 0.25\textdegree\ cells,
  with USGS Quaternary faults and 1351 ComCat epicentres of magnitude 4 and above
  shallower than 40\,km, 1900--2025. (b--d)~Length-weighted strikes of (b) mapped and
  (c) Quaternary faults and (d) nodal-plane strikes of 103 shallow Global CMT
  mechanisms, in 10\textdegree\ sectors normalised to the largest sector.
  (e)~Cumulative distribution of the distance from each epicentre to the nearest
  Quaternary fault, compared with 6755 uniformly random points inside Afghanistan.
  Software used for plotting figure: GMT 6.6.0. Data: USGS faults of Afghanistan
  \cite{Doebrich2006}; USGS Quaternary faults \cite{Ruleman2007}; USGS ComCat; Global
  CMT \cite{Ekstrom2012}. Source: authors.}\label{fig:faultdens}
\end{figure}

\subsection{Recurrence and Sensitivity Analysis}
The Gutenberg--Richter parameters give return periods of $M_w \geq 7$ events of 11
years for HKP (95\,\% interval 7--19), 19 years for WEST (11--47), 70 years for NEC
(31--269) and 87 years for PBZ (31--355) (\cref{fig:sensitivity}b). The observed
numbers of such events in 1900--2025 correspond to recurrence intervals of 10.5
years in HKP, 21 years in NEC and PBZ, and 25 years in WEST. The probability of at
least one $M_w \geq 7$ event in 50 years is 0.99 for HKP, 0.93 for WEST (0.65--0.99),
0.51 for NEC (0.17--0.81) and 0.44 for PBZ (0.13--0.80).

\begin{figure}[htbp]\centering
  	\includegraphics[width=0.98\textwidth]{fig_sensitivity}
  \caption{Sensitivity of the b-values and recurrence of large earthquakes. (a)
  b-values with Shi--Bolt standard errors for each domain under alternative
  choices. These are the reference nearest-neighbour catalogue with $M_c = 5.5$, the
  undeclustered and Gardner--Knopoff catalogues and the catalogue without default
  depths. Further rows use completeness magnitudes of 5.3 and 5.8 and the
  converted catalogue with the regional and Scordilis conversions. (b)
  Return periods of $M_w \geq 7$ events from the Gutenberg--Richter parameters, with
  bootstrap 95\,\% intervals (filled bars), and from the observed counts of
  1900--2025 (hatched bars), on a logarithmic axis. Software used for plotting
  figure: Python 3.14, Matplotlib, NumPy. Data: USGS ComCat; ISC-GEM. Source:
  authors.}
  \label{fig:sensitivity}
\end{figure}

The ranking of the domains is stable under the declustering and depth choices
(\cref{fig:sensitivity}a). With Gardner--Knopoff declustering, the b-values decrease
by 0.06--0.27 in every domain, and WEST remains the lowest at 0.64. Raising $M_c$ to
5.8 changes the values by less than 0.13 where enough events remain. Lowering $M_c$
to 5.3 raises all values by 0.13--0.41. Below 5.5, converted magnitudes enter the
sample. The converted catalogue yields systematically higher b-values: 1.27--1.85
with the Scordilis conversion and 1.03--1.85 with the regional conversion, at
completeness magnitudes of 4.7--5.0.

Taken together, the four domains differ consistently in b-value, depth, faulting
style, strain and the timing of moment release, and their ranking holds under every
declustering, completeness and depth choice tested. The Discussion interprets these
contrasts in terms of the processes that produce them, examines why the shallow
domains host the destructive earthquakes and considers what the results imply for
disaster risk reduction in Afghanistan.

% ===========================================================================
\section{Discussion}\label{sec:discussion}

\subsection{Three Seismicity Regimes and Their Tectonic Origin}
The results define four seismotectonic domains. They group into three regimes:
intermediate-depth seismicity beneath the Hindu Kush and Pamir, crustal deformation
along the plate-boundary zone and the north-eastern ranges, and sparse but
strike-slip-dominated seismicity in the western interior. \Cref{fig:synthesis}
summarises their properties.
The contrasts in b-value follow the state of stress and the structure of each
domain. Laboratory experiments and global compilations show that b decreases as
differential stress increases, so that thrust faulting yields lower b-values than
strike-slip and normal faulting \cite{Scholz1968, Schorlemmer2005, Scholz2015}. The
intermediate-depth zone fits this relation. Its b-value of 0.97 accompanies a
dominance of thrust-class mechanisms and a tightly clustered N--S shortening axis.
In a near-vertical slab, such mechanisms have steeply plunging T axes and therefore
record down-dip stretching of the lithosphere that sinks beneath the Hindu Kush
\cite{Zhan2016, Molnar2019, Kufner2017}. Five cells of the nest carry 96\,\% of the
intermediate-depth strain rate, which shows that this stretching is focused in a
small volume, as expected for the necking of a detaching slab \cite{Lister2008,
Kufner2016, Kufner2017}. The two depth modes near 115 and 204\,km and the lowest
gridded b-values beneath the southern tip of the slab further resolve vertical and
lateral variations of stress within this volume \cite{Pegler1998, Sippl2013}.

The north-eastern crust and the plate-boundary zone have the highest b-values,
1.33--1.38. Both domains are structurally complex. The north-eastern crust combines
thrusting with the largest share of normal faulting, and the plate-boundary zone
combines thrusting on the Sulaiman and Himalayan fronts with strike-slip faulting
along the Chaman system. Heterogeneous fault populations of this kind, together with
the partly creeping segments of the Chaman fault \cite{Fattahi2016, Barnhart2017,
Dalaison2021}, release strain in many moderate events relative to large ones, which
raises b. The Kostrov field supports this interpretation. Shallow strain rates span
four orders of magnitude and are highest along the margins of Afghanistan, where
geodetic networks record the partitioning of India--Eurasia convergence between
strike-slip and thrust structures \cite{Mohadjer2010, Szeliga2012}.

The western interior has the lowest b-value, 0.69, although strike-slip mechanisms
dominate it. Here the low value reflects few moderate events relative to the large
earthquakes of 1997 Zirkuh and 2013 Saravan rather than high stress alone.
Deformation in this domain is slow and distributed over long faults such as the
Herat fault \cite{Wheeler2005, Ruleman2007, Walker2004}, so the instrumental window
samples only part of its seismic cycle. The NE--SW P axis of the western interior,
distinct from the N--S axes of the other domains, records the transfer of
Arabia--Eurasia convergence into eastern Iran and western Afghanistan
\cite{Vernant2004, Walker2004}. Seismicity in Afghanistan therefore reflects three
processes: stretching of a sinking slab, partitioned shortening along the
India--Eurasia margin and slow distributed deformation of the Afghan interior.

\subsection{Destructive Crustal Earthquakes Outside the Hindu Kush Nest}
All destructive earthquakes of 2022--2025 occurred at 4--28\,km depth in the
crustal domains. The 2022 Paktika and 2025 Kunar events lie in the plate-boundary
zone, the 2023 Herat sequence in the western interior and the 2025 Kholm event in
the north-eastern crust (\cref{fig:synthesis}). Together they released only 0.13\,\%
of the seismic moment of the study period. Their impact therefore follows from three
factors that catalogue statistics do not capture. First, shallow hypocentres beneath
inhabited valleys concentrate strong shaking in small epicentral areas
\cite{Jahed2024}. Second, the adobe and stone dwellings that dominate rural housing
collapse under moderate ground motion \cite{Akhundzadah2025, WorldBank2024}. Third,
some sources were unknown before they ruptured. The Herat sequence broke a set of
north-dipping blind thrusts north-west of the city of Herat in five mainshocks
within nine days \cite{Zhao2025, Fu2025, Xu2025, Zhong2026, Toker2026}. Faults of
this kind do not reach the surface and are absent or incomplete in the Quaternary
fault maps.

The fault test shows why this matters. Instrumental epicentres lie a median of
3.5\,km from mapped Quaternary faults, so the mapped network captures most of the
recorded shallow activity. The same analysis shows that the Helmand and Farah blocks,
with few mapped faults, also record few earthquakes. A short catalogue and an
incomplete fault map can thus agree on low activity where buried structures
accumulate strain slowly. Experience in other slowly deforming continental regions
shows that such areas host infrequent but damaging earthquakes, which hazard maps
built on short catalogues tend to underestimate \cite{Stein2012, Liu2016, Calais2016,
England2011}. In the western interior, the modelled return period of $M_w \geq 7$
events is about two decades. Large shallow earthquakes are therefore a recurrent
part of its behaviour, even allowing for the weight of the 2013 Saravan earthquake
in its statistics. The destructive events of 2022--2025 are therefore not anomalies
but a sample of the activity that the crustal domains sustain over longer periods.

\begin{figure}[htbp]\centering
  \includegraphics[width=0.74\textwidth]{fig_synthesis}
  \caption{Synthesis of the seismicity regimes of Afghanistan. (a)~Crustal domains,
  shaded by the majority domain of the events in each 0.5\textdegree\ cell, and the
  footprint of the intermediate-depth domain (blue outline). Beachballs show the
  summed Global CMT moment tensor of each domain, crosses the principal horizontal
  strain axes from Kostrov summation and stars the destructive earthquakes of
  2022--2025. (b)~Key statistics of each domain: b-value for $M_w \geq 5.5$ in
  1964--2025, dominant faulting style and mean P axis, share of the seismic moment
  released in 1900--2025, median depth and modelled return period of $M_w \geq 7$
  events. Software used for plotting figure: GMT 6.6.0. Data: this study (merged
  ComCat and ISC-GEM catalogue; Global CMT); SRTM15+. Source: authors.}
  \label{fig:synthesis}
\end{figure}

\subsection{Comparison with Previous Hazard and Seismicity Studies}
The domain b-values agree with earlier regional estimates where the data and
methods are comparable (\cref{tab:comparison}). Drakopoulos and Srivastava
\cite{Drakopoulos1972} obtained $b = 1.5 \pm 0.2$ for shallow and $1.2 \pm 0.2$ for
intermediate-depth Hindu Kush earthquakes of 1963--1971 and found no decrease of b
with depth. Rehman et al.\ \cite{Rehman2017} identified three areas with b below 0.9
within the intermediate-depth zone. Rehman et al.\ \cite{Rehman2023} later reported
b-values of 0.99--1.47 across the Pamir--Hindu Kush, with the lowest values in the
Hindu Kush and the northern Pamir. The intermediate-depth value of 0.97 found here
lies at the lower end of this range, and the low gridded values beneath the southern
tip of the slab coincide with the low-b intermediate-depth areas mapped earlier. The
higher values of the north-eastern crust and the plate-boundary zone agree with the
higher shallow b-values of the earlier studies. The remaining differences in absolute
value reflect magnitude type and threshold. The earlier studies used body-wave or
mixed magnitudes and lower completeness magnitudes, and the sensitivity tests of
this study show that both choices raise b (\cref{fig:sensitivity}a).

The national hazard map of Afghanistan combined a catalogue of earthquakes of
magnitude 4.5 and above since 1964, with depths to 250\,km, and the Quaternary fault
database \cite{Boyd2007, Ruleman2007}. The present analysis complements this model
with objectively delineated domains and with domain-specific b-values, depth
distributions, faulting styles and strain directions. These parameters can define
separate crustal and intermediate-depth source zones in an updated hazard model. They also
identify the western interior as a domain whose low b-value implies a larger share
of large earthquakes than a regional average would suggest \cite{Giardini1999,
Silva2020}.

\begin{table}[htbp]\centering
  \caption{Comparison of b-values from this study with published estimates for the
  Hindu Kush--Pamir region. NND, nearest-neighbour declustering; USCGS, U.S. Coast
  and Geodetic Survey; NEIC, USGS National Earthquake Information Center.}
  \label{tab:comparison}
  \small
  \begin{tabular}{>{\raggedright\arraybackslash}p{3.1cm}>{\raggedright\arraybackslash}p{3.4cm}>{\raggedright\arraybackslash}p{1.7cm}>{\raggedright\arraybackslash}p{2.0cm}>{\raggedright\arraybackslash}p{3.8cm}}
    \toprule
    Study & Region and depth & Period & $b$ & Data and method \\
    \midrule
    Drakopoulos and Srivastava \cite{Drakopoulos1972} & Hindu Kush, 0--100\,km
      & 1963--1971 & $1.5 \pm 0.2$ & USCGS magnitudes; least squares \\
     & Hindu Kush, 100--200\,km & 1963--1971 & $1.2 \pm 0.2$ & as above \\
    Rehman et al.\ \cite{Rehman2017} & Hindu Kush, 70--300\,km & to 2007
      & $< 0.9$ in three areas & b-value maps of a declustered catalogue \\
    Rehman et al.\ \cite{Rehman2023} & Pamir--Hindu Kush, four zones & 1949--2018
      & 0.99--1.47 & NEIC, $M \geq 3.0$, declustered; lowest b in Hindu Kush \\
    This study, HKP & Hindu Kush--Pamir, $\geq 70$\,km & 1964--2025
      & $0.97 \pm 0.07$ & $M_w \geq 5.5$, NND; Aki--Utsu, Shi--Bolt \\
    This study, NEC & North-eastern crust & 1964--2025 & $1.33 \pm 0.16$ & as above \\
    This study, PBZ & Plate-boundary zone & 1964--2025 & $1.38 \pm 0.19$ & as above \\
    This study, WEST & Western interior & 1964--2025 & $0.69 \pm 0.11$ & as above \\
    \bottomrule
  \end{tabular}
\end{table}

\subsection{Implications for Disaster Risk Reduction in Afghanistan}
The domain statistics translate into four priorities for the Afghanistan National
Disaster Management Authority (ANDMA) and its partners. First, hazard zoning benefits
from treating the shallow crustal domains separately from the Hindu Kush nest. The
nest dominates the catalogue and the moment budget, but its depths of 150--250\,km
spread and attenuate shaking. The crustal domains, in contrast, produced every
lethal earthquake of 2022--2025. Separate source zones with their own b-values, depth
distributions and faulting styles keep crustal hazard from being diluted by deep
seismicity \cite{Boyd2007, Silva2020}. Second, the western interior and the Herat
region merit high priority. Their low b-value and the blind-thrust ruptures of 2023
show that large shallow earthquakes recur there although moderate events are rare.

Third, building-code provisions and retrofitting of adobe and stone dwellings are
most effective where shallow sources lie beneath dense rural settlement. Such areas
include the Herat valley, the Kabul--Jalalabad corridor, Paktika and Khost, and the
Kunar valley \cite{Akhundzadah2025, Ismail2016}. Fourth, preparedness planning should
anticipate earthquake sequences. The Herat sequence produced four $M_w$ 6.3 events in
nine days, so that shelters, damage inspection and search-and-rescue procedures must
cope with repeated strong shaking of already weakened buildings \cite{WorldBank2024,
Zhao2025}. Because the catalogue, scripts and maps are open, ANDMA staff and
partner institutions can update these statistics as new earthquakes occur.

\subsection{Limitations}
The domains and their ranking are robust, but five aspects of the data bound the
precision of the results. First, the local seismic network of Afghanistan is sparse,
so ComCat epicentres carry errors of 10--20\,km and many crustal depths are fixed at
default values. These errors blur the gridded statistics and the fault-distance test
at scales below about 20\,km, and default-depth events remain in the crustal b-value
grid. Second, the domain b-values rest on events of $M_w$ 5.5 and above since 1964,
which limits sample sizes in the crustal domains; the gridded body-wave b-values are
therefore read as relative contrasts. Third, converted magnitudes raise b, as the
sensitivity tests show, so absolute values from converted catalogues are not
directly comparable with moment-magnitude estimates. Fourth, the Global CMT
catalogue is complete only above about $M_w$ 5, so Kostrov summation leaves the
interior between 62\textdegree E and 66\textdegree E unresolved and captures only
part of the deformation of slowly deforming regions. Fifth, the fault database is
generalised at 1:850\,000 scale and omits blind structures such as those that
ruptured near Herat. In addition, the western domain extends into eastern Iran, and
its moment budget and recurrence are influenced by the 2013 Saravan earthquake, an
event at about 80\,km depth beneath the Iranian side of the border.

\subsection{Future Directions}
Three lines of work follow from these results. First, a companion study will
analyse the subsurface structure of the northern margin of the Afghan platform,
where the 2025 Kholm earthquake occurred at 28\,km depth, using gravity, magnetic and
topographic data to locate buried thrusts and basement faults. Second, city-scale
scenario shaking for Herat, Kabul and Jalalabad will combine the domain parameters
of this study with ground-motion models and building inventories to estimate damage
for representative earthquakes. Third, denser broadband and strong-motion networks,
operated jointly with ANDMA and regional partners, would lower completeness
magnitudes, constrain crustal depths and allow b-values to be mapped at the scale of
individual fault systems. Relocation of the instrumental catalogue with regional
velocity models would sharpen the fault-distance test and the depth sections.

% ===========================================================================
\section{Conclusions}\label{sec:conclusions}
This study compared the statistical seismicity regimes of Afghanistan using open
data and a single reproducible workflow. Unsupervised clustering of a merged and
declustered catalogue resolves four seismotectonic domains: the Hindu Kush--Pamir
intermediate-depth zone, the north-eastern crust, the plate-boundary zone and the
western interior. These domains differ systematically in b-value, depth, faulting
style, strain and the timing of moment release.

The intermediate-depth zone dominates the catalogue and the moment budget and
records the stretching of a near-vertical slab under N--S shortening. The
north-eastern crust and the plate-boundary zone release strain in many moderate
events across heterogeneous fault systems, whereas the western interior has the
lowest b-value and a NE--SW shortening direction. Shallow seismicity concentrates
near the mapped Quaternary faults. The destructive earthquakes of recent years all
occurred in the crust, several of them on blind thrusts, while releasing only a small
fraction of the seismic moment.

The study provides the first nationwide, domain-resolved statistical regionalisation
of Afghan seismicity, together with open scripts and maps that partner institutions
can reproduce. For disaster risk reduction, it supports hazard zoning that separates
shallow crustal sources from the deep nest and gives priority to the western
interior and to populated valleys above shallow thrusts. The main limitation is the
reliance on teleseismic locations and magnitudes. Future work will image the buried
faults of the northern platform margin and model scenario shaking for the major
cities of Afghanistan.

% ===========================================================================
\section*{Data and Code Availability}
All data used in this study are openly available: USGS ComCat
(\url{https://earthquake.usgs.gov/earthquakes/search/}), ISC-GEM and the ISC
Bulletin (\url{https://www.isc.ac.uk}), Global CMT (\url{https://www.globalcmt.org}),
USGS Open-File Reports 2006-1038 and 2007-1103, PB2002, SRTM15+ and GSHHG.
The Python and GMT scripts used for catalogue processing, statistical analysis and
mapping are openly available on GitHub
(\url{https://github.com/paulinelemenkova/seismotectonic-regime-partitioning}) and
archived on Zenodo (\url{https://doi.org/10.5281/zenodo.24098106}).

\section*{Funding}
This research was supported by T\"{U}B\.{I}TAK --- B\.{I}DEB 2221 (Science
Fellowships and Grant Programmes), grant number 2221, reference number
B.14.2.TBT.0.06.01.02-220-859260.

\section*{CRediT Author Statement}
\textbf{Polina Lemenkova}: Conceptualization, Methodology, Software, Data curation,
Formal analysis, Validation, Investigation, Visualization, Writing -- original draft,
Writing -- review \& editing. \textbf{Abdullah Can Z\"{u}lfikar}: Conceptualization,
Supervision, Project administration, Funding acquisition, Writing -- review \&
editing.

\section*{Acknowledgments}
The authors thank the U.S. Geological Survey, the International Seismological
Centre, the Global CMT Project and the NOAA National Centers for Environmental
Information for providing open earthquake and fault data. They also thank the developers of
GMT, Python, NumPy, SciPy, pandas, scikit-learn and Matplotlib. This work contributes to
the ITU--T\.{I}KA--JICA training programme on earthquake hazard and disaster risk
for staff of the Afghanistan National Disaster Management Authority.

\section*{Conflict of Interest Statement}
The authors declare no competing interests.

% ===========================================================================
\appendix
\section{Code Listings}\label{secA1}

The listings below give the core steps of the catalogue analysis and the map
scripts; the complete scripts are available as described under Data and Code
Availability.

\begin{lstlisting}[style=sh,caption={GMT script for the depth-coloured
seismicity map with depth and magnitude histograms (GMT 6.6.0; excerpt of
fig\_catalogue\_overview.sh).},label={lst:seismap}]
gmt grdcut @earth_relief_30s -R$R -Grelief.nc
gmt grdgradient relief.nc -A315 -Nt0.8 -Gint.nc
gmt makecpt -Cgray -T-14000/6000/250 > relief.cpt
gmt makecpt -Cturbo -T0/300/10 -I -D > depth.cpt
gmt begin fig_catalogue_overview pdf,png E300
gmt grdimage relief.nc -Iint.nc -Crelief.cpt -R$R -J$J
gmt coast -Di -W0.3p,gray30 -N1/0.6p,gray25,- -I1/0.35p,steelblue \
          -C+l+gsteelblue -A50 -Bxa2f1g2 -Bya2f1g2 -BWSen
gmt plot qf.gmt -W0.3p,gray20                     # Quaternary faults
gmt plot af.txt -W1.4p,black                      # national border
# colour = depth, symbol size grows with magnitude
awk '{printf "%s %s %s %.3f\n",$1,$2,$3,0.03*exp(($4-3)*log(1.9))}' eq.txt \
  | gmt plot -Sc -Cdepth.cpt -W0.1p,black@40
awk '{print $1,$2}' key.txt | gmt plot -Sa0.55c -W1.2p,black -Gwhite@100
gmt colorbar -Cdepth.cpt -DJBC+w17c/0.32c+o0/1.0c+h+ef \
             -Bxa50f10+l"Hypocentral depth (km)"
gmt histogram depth.txt -R0/300/1/10000 -JX7.9c/3.4cl -T10 -Gsteelblue3 \
   -W0.3p,white -Bxa50f10+l"Depth (km)" -Bya1pf3+l"Events" -BWSne -Y-6.0c
gmt histogram mag.txt -R3.9/8/1/10000 -JX7.9c/3.4cl -T0.1 -Gfirebrick3 \
   -W0.3p,white -Bxa0.5f0.1+l"Magnitude" -Bya1pf3+l"Events" -BWSne -X9.1c
gmt end
\end{lstlisting}

\begin{lstlisting}[style=sh,caption={GMT projection of hypocentres and topography
onto the section lines and plotting of one depth section (GMT 6.6.0; excerpt of
fig\_depth\_sections.sh).},label={lst:sections}]
while read x0 y0 x1 y1 L; do                  # one line per section
  gmt project eq.txt -C$x0/$y0 -E$x1/$y1 -W-$HW/$HW -Lw -Q -Fpz > sec_$L.txt
  gmt project -C$x0/$y0 -E$x1/$y1 -G1 -Q > pts_$L.txt
  gmt grdtrack pts_$L.txt -Grelief.nc | awk '{print $3, $4/1000}' > topo_$L.txt
done < lines.txt
msize() { awk '{s=0.035*1.55^($3-3); print $1, $2, $2, s}'; }
sec() {  # L  zmax  height(cm)  yshift
  local L=$1 ZM=$2 H=$3 YS=$4
  local LEN=$(awk -v l=$L '$1==l{print $2}' len.txt)
  gmt basemap -R0/$LEN/-1/7 -JX16.2c/1.0c -Y$YS -Bya3f1+l"km" -BWsen
  awk 'NR==1{print $1,-1} {print} END{print $1,-1}' topo_$L.txt \
    | gmt plot -G150/150/150 -W0.4p,gray20
  gmt basemap -R0/$LEN/0/$ZM -JX16.2c/-${H}c -Y-${H}c \
    -Bxa100f25+l"Distance along section (km)" -Bya50f25g50+l"Depth (km)" -BWSen
  msize < sec_$L.txt | gmt plot -Sc -Cdep.cpt -W0.15p,gray20
}
sec A 300 4.6 -2.5c
\end{lstlisting}

\begin{lstlisting}[caption={Frohlich classification of Global CMT mechanisms read
from NDK records and GMT plotting of the mechanisms coloured by faulting style
(Python 3 and GMT 6.6.0; excerpt of fig\_focal\_mechanisms.sh).},label={lst:meca}]
for i in range(0, len(lines) - 4, 5):              # five lines per NDK event
    l1, l2, l3, l4, l5 = lines[i:i + 5]
    c = l3.split(); lat, lon, dep = float(c[3]), float(c[5]), float(c[7])
    m = l4.split(); ex = int(m[0])
    mt = [float(m[k]) for k in (1, 3, 5, 7, 9, 11)]
    v = l5.split()
    tpl, bpl, ppl = float(v[2]), float(v[5]), float(v[8])   # axis plunges
    if tpl >= 50:   k = "TF"                       # thrust
    elif ppl >= 60: k = "NF"                       # normal
    elif bpl >= 60: k = "SS"                       # strike-slip
    else:           k = "OB"                       # oblique
    rows.append(dict(lon=lon, lat=lat, dep=dep, mt=mt, ex=ex, k=k))
# --- GMT (shell): deep events grey beneath, shallow events by class
# gmt meca meca_in_all.txt -Sm0.22c -Ggray30 -Ewhite -W0.15p,gray10
# for k in TF NF SS OB; do
#   eval col=\$C_$k
#   gmt meca meca_sh_$k.txt -Sm0.38c -G$col -Ewhite -W0.25p,black
# done
\end{lstlisting}

\begin{lstlisting}[caption={Nearest-neighbour declustering with a Gaussian-mixture
threshold, HDBSCAN clustering of background hypocentres, grouping of clusters into
domains and k-nearest-neighbour labelling of all events (Python 3, scikit-learn;
excerpt of analysis.py).},label={lst:clustering}]
v = le[np.isfinite(le)].reshape(-1, 1)            # log10 eta of each event
gm = GaussianMixture(2, random_state=0).fit(v)
mu = np.sort(gm.means_.ravel())
grid = np.linspace(mu[0], mu[1], 4000)
lo = int(np.argmin(gm.means_.ravel()))
pp = gm.predict_proba(grid.reshape(-1, 1))[:, lo]
thr = float(grid[np.argmin(np.abs(pp - 0.5))])    # equal membership
cat["nnd_bg"] = ~(le < thr)                       # background events

def feats(d, w=3.0):                  # easting, northing, 3 x depth (km)
    return np.c_[(d.longitude - 66.0) * 111.32 * np.cos(np.radians(33.0)),
                 (d.latitude - 33.0) * 110.57, d.depth * w]

bg = cat[cat.nnd_bg]
lab = HDBSCAN(min_cluster_size=80, min_samples=20).fit_predict(feats(bg))

def group(c):                                     # rules on cluster medians
    if c["depth_med"] >= 70:
        return "HKP"
    if c["lon"] < 64.5:
        return "WEST"
    if c["lat"] < 35.2 or c["lon"] > 72.5 and c["lat"] < 36.0:
        return "PBZ"
    return "NEC"

lab_dom = np.array([gmap.get(L, "NOISE") for L in lab])
ok = lab_dom != "NOISE"
knn = KNeighborsClassifier(n_neighbors=15).fit(feats(bg)[ok], lab_dom[ok])
cat["domain"] = knn.predict(feats(cat))
\end{lstlisting}

\begin{lstlisting}[caption={Fault-density gridding in 0.25\textdegree\ cells and
comparison of epicentre-to-fault distances with uniformly random points by a
two-sample Kolmogorov--Smirnov test (Python 3, NumPy, SciPy; excerpt of
fig\_fault\_density.sh).},label={lst:faults}]
F = pieces(read_lines("flt.gmt"))                 # midpoint lon, lat, length, azimuth
d = 0.25; xs = np.arange(60, 75.5, d); ys = np.arange(29, 38.75, d)
G = np.zeros((len(ys), len(xs)))
ix = np.floor((F[:, 0] - 60) / d).astype(int)
iy = np.floor((F[:, 1] - 29) / d).astype(int)
ok = (ix >= 0) & (ix < len(xs)) & (iy >= 0) & (iy < len(ys))
np.add.at(G, (iy[ok], ix[ok]), F[ok, 2])         # summed length per cell
area = R**2 * np.radians(d) * (np.sin(np.radians(ys + d)) - np.sin(np.radians(ys)))
G = G / area[:, None]                             # km per km^2

def proj(lon, lat):                               # local equirectangular (km)
    return np.c_[np.radians(lon - 67.5) * R * np.cos(np.radians(lat)),
                 np.radians(lat) * R]

vx = []                                           # densify faults to <= 0.5 km
for sgm in read_lines("qf.gmt"):
    for (x1, y1), (x2, y2) in zip(sgm[:-1], sgm[1:]):
        L = np.hypot((x2 - x1) * 111.2 * np.cos(np.radians(y1)), (y2 - y1) * 111.2)
        t = np.linspace(0, 1, max(2, int(np.ceil(L / 0.5)) + 1))
        vx.append(np.c_[x1 + t * (x2 - x1), y1 + t * (y2 - y1)])
vx = np.vstack(vx); tree = cKDTree(proj(vx[:, 0], vx[:, 1]))
de = tree.query(proj(pe[:, 0], pe[:, 1]))[0]      # epicentres
dr = tree.query(proj(pr[:, 0], pr[:, 1]))[0]      # random points
ks = stats.ks_2samp(de, dr)
\end{lstlisting}

\bibliographystyle{plainnat}
\bibliography{Bibliography}

\end{document}
